\chapter[Band 16: Wohlordnungen und Auswahlaxiom]{Band 16 auf einen Blick}
\noindent\textbf{Wohlordnungen und Auswahlaxiom}\par\medskip
\noindent\emph{Lesart.} Eine Wohlordnung ist eine totale Ordnung,
in der jede nichtleere Teilmenge ein kleinstes Element hat.
\emph{Wichtig für den Stand dieses Bandes:} Zermelos Konstruktionsprinzip
ist im Original als Axiom angesetzt. Die Richtung vom Auswahlaxiom zum
Wohlordnungssatz wird unter Verwendung dieses zusätzlichen Prinzips
hergeleitet; die Übersicht gibt es nicht als dort bewiesenen Rekursionssatz aus.
\begin{BandUebersicht}
\textbf{Wohlordnung}
Zusätzlich zu \(\TotOrd{A,\leq}\) wird gefordert:
\UebFormel{\begin{aligned}
&T\subseteq A\dsep T\neq\varnothing\\
&\hspace{10mm}\vdash\exists m\in T\,\forall x\in T\,(m\leq x).
\end{aligned}}
\UebQuellen{\FormulaRefAuto{Wohlordnung}[def];
\FormulaRefAuto{WellOrderPrinciple}[ax]}
&
\textit{Natürliches Beispiel}
\UebFormel{\WellOrd{\mathbb N,\leq_{\mathbb N}}.}
Der Satz verbindet die totale Ordnung aus Band 15 mit dem
Wohlordnungsprinzip aus Band 10.
\UebQuellen{\FormulaRefAuto{NaturalNumbersWellOrder}[thm]}\\
\addlinespace[.8ex]
\textbf{Wohlordenbarkeit und Wohlordnungssatz}
\UebFormel{\begin{aligned}
&\WellOrdable{A}\coloneqq\\
&\qquad\exists\preccurlyeq\;\WellOrd{A,\preccurlyeq},\\
&\WOP\coloneqq\forall A\,\WellOrdable{A}.
\end{aligned}}
\UebQuellen{\FormulaRefAuto{WellOrdable}[def];
\FormulaRefAuto{WOP}[def]}
&
\textit{Wohlordnungssatz liefert Auswahl}
\UebFormel{\WOP\vdash\ACsur.}
Die Richtung wird über die kleinsten Elemente der Fasern einer
Surjektion bewiesen.
\UebQuellen{\FormulaRefAuto{WOPToACsur}[thm]}\\
\addlinespace[.8ex]
\textbf{Auswahlrelation auf nichtleeren Teilmengen}
\UebFormel{\begin{aligned}
&\ChoiceRel{A}\coloneqq\\
&\qquad\{(X,a)\in\PowNE{A}\times A\\
&\hspace{20mm}\mid a\in X\}.
\end{aligned}}
\UebQuellen{\FormulaRefAuto{ChoiceRelation}[def]}
&
\textit{Die Auswahlrelation ist total}
\UebFormel{\begin{aligned}
&\TotRel{\ChoiceRel{A},\PowNE{A},A},\\
&(X,a)\in\ChoiceRel{A}\vdash a\in X.
\end{aligned}}
\UebQuellen{\FormulaRefAuto{ChoiceRelTotal}[thm];
\FormulaRefAuto{ChoiceRelElim}[thm]}\\
\addlinespace[.8ex]
\textbf{Auswahlfunktion auf der nichtleeren Potenzmenge}
\UebFormel{\begin{aligned}
&\ChoicePow{C}{A}\coloneqq\\
&\quad C\colon\PowNE{A}\to A\\
&\quad\land\forall X\in\PowNE{A}\,(C(X)\in X).
\end{aligned}}
\UebQuellen{\FormulaRefAuto{ChoicePow}[def]}
&
\textit{AC liefert eine solche Auswahlfunktion}
\UebFormel{\ACsur\vdash\exists C\,\ChoicePow{C}{A}.}
\UebQuellen{\FormulaRefAuto{ACsurChoicePow}[thm]}\\
\addlinespace[.8ex]
\textbf{Zermelos Konstruktionsprinzip -- Axiom im Original}
\UebFormel{\begin{aligned}
&\ChoicePow{C}{A}\\
&\hspace{10mm}\vdash\exists\preccurlyeq\;\WellOrd{A,\preccurlyeq}.
\end{aligned}}
Dieses Prinzip steht auf der Axiomseite des gegenwärtigen Aufbaus.
\UebQuellen{\FormulaRefAuto{ZermeloConstruction}[ax]}
&
\textit{Unter diesem Prinzip: Auswahl impliziert Wohlordnung}
\UebFormel{\begin{aligned}
&\exists C\,\ChoicePow{C}{A}\vdash\WellOrdable{A},\\
&\ACsur\vdash\WOP,\\
&\ACsur\eqvdash\WOP.
\end{aligned}}
Die letzte Zeile kombiniert diese Richtung mit dem unabhängig
von der Konstruktion bewiesenen Rückweg.
\UebQuellen{\FormulaRefAuto{ZermeloWO}[thm];
\FormulaRefAuto{ACsurToWOP}[thm];
\FormulaRefAuto{ACsurEquivWOP}[thm]}\\
\addlinespace[.8ex]
\textbf{Faserminimenge}
Für \(\WellOrd{B,\preccurlyeq}\), \(G\colon B\sur A\):
\UebFormel{\begin{aligned}
&L_{G,\preccurlyeq}\coloneqq\{b\in B\mid\forall c\in B\\
&\qquad(G(c)=G(b)\rightarrow b\preccurlyeq c)\}.
\end{aligned}}
\UebQuellen{\FormulaRefAuto{FiberMinSet}[def]}
&
\textit{Faserminima liefern eine Sektion}
\UebFormel{\begin{aligned}
&\forall X\in\Fib_G[A]\\
&\qquad\exists!b\,(b\in X\land b\in L_{G,\preccurlyeq}),\\
&\exists F\colon A\to B\;G\circ F=\Id_A.
\end{aligned}}
\UebQuellen{\FormulaRefAuto{WellOrdFiberChoiceSet}[thm];
\FormulaRefAuto{WellOrdFibersSection}[thm]}\\
\end{BandUebersicht}
